% ============================================================
%  Chapter 1 -- The Extended Real Number System, 1.23
% ============================================================
\rsection{The Extended Real Number System}

\begin{signpost}[A convenience, not a number system]
One definition, no theorems. Rudin is buying a notational convenience: he wants
to write $\sup E = +\infty$ for unbounded $E$, and $\lim s_n = +\infty$ in
Chapter 3, without special-casing every statement. Read it as an agreement about
symbols.

\smallskip
The thing to carry away is what he pointedly does \emph{not} do. The extended
reals are given an order but not an arithmetic --- ``does not form a field'' is
stated flatly --- and the conventions in (a)--(c) leave $\infty - \infty$,
$0 \cdot \infty$, and $\infty/\infty$ undefined. Those gaps are deliberate and
permanent; they are why $\liminf$ and $\limsup$ in 3.16 need care, and why
Theorem 3.42 has the hypotheses it does.
\end{signpost}

\ritem{1.23}{Definition}
The \emph{extended real number system} consists of the real field $\R$ and two
symbols, $+\infty$ and $-\infty$. We preserve the original order in $\R$ and
define
\[ -\infty < x < +\infty \]
for every $x \in \R$.

It is then clear that $+\infty$ is an upper bound of every subset of the
extended real number system, and that every nonempty subset has a least upper
bound.\kn{This is the payoff. Of the two hypotheses in Definition 1.10, exactly
one becomes vacuous: ``bounded above'' is now automatic, since $+\infty$ bounds
every subset. ``Nonempty'' is still required, and Rudin keeps it in the sentence
alongside. What has been bought, at the cost of field structure, is a $\sup$
that exists for every nonempty set whether or not it is bounded.} If,
for example, $E$ is a nonempty subset of real numbers which is not bounded above
in $\R$, then $\sup E = +\infty$ in the extended real number system.

Exactly the same remarks apply to lower bounds.

The extended real number system does not form a field, but it is customary to
make the following conventions:
\begin{enumerate}[label=(\alph*), leftmargin=2.2em, itemsep=3pt, topsep=4pt]
  \item If $x$ is real then
        \[ x + \infty = +\infty, \qquad x - \infty = -\infty, \qquad
           \frac{x}{+\infty} = \frac{x}{-\infty} = 0. \]
  \item If $x > 0$ then $x \cdot (+\infty) = +\infty$,
        $x \cdot (-\infty) = -\infty$.
  \item If $x < 0$ then $x \cdot (+\infty) = -\infty$,
        $x \cdot (-\infty) = +\infty$.
\end{enumerate}

When it is desired to make the distinction between real numbers on the one hand
and the symbols $+\infty$ and $-\infty$ on the other quite explicit, the former
are called \emph{finite}.

\begin{pitfall}[Read the conventions as a closed list]
Every case Rudin lists involves one infinity and one \emph{finite} $x$. Anything
outside that pattern is simply undefined --- among them $(+\infty)+(-\infty)$,
$0\cdot(\pm\infty)$, every quotient of two infinities, and every sum or product
of two infinite symbols not covered above.

The omissions are deliberate, not oversights, because no assignment preserves
the ordinary rules: if $\infty - \infty$ were given a value $c$, then adding $1$
to the first term would force $c = c+1$. (Division by zero is a separate matter
--- $x/0$ is undefined in any field, by (M5), and has nothing to do with the
extended system.)

The practical consequence: an expression like $\sup E + \sup F$ is not
automatically meaningful once $\sup$ may be infinite, and theorems in Chapter 3
that add limits carry hypotheses excluding these cases. When a later proof says
``provided the right-hand side is not of the form $\infty - \infty$,'' this is
what it is protecting against.
\end{pitfall}

\begin{deeper}[Why not just add $\infty$ and keep the field axioms?]
Because it is impossible, and 1.16(b) shows why in two lines. Suppose
$+\infty$ had a multiplicative inverse; convention (a) already forces
$x/(+\infty) = 0$, so $1/(+\infty) = 0$, and then
$1 = (+\infty) \cdot 0 = 0$ by 1.16(a) --- contradicting $1 \ne 0$ in (M4).

So the extended reals are an ordered set with a partial arithmetic, and that is
all they can be. Rudin's phrasing --- ``it is customary to make the following
conventions'' --- is exact: these are conventions, not theorems, and they are
adopted because they make limit statements uniform, not because they are forced.
\end{deeper}
